\documentclass[10pt,a4paper]{amsart}
\usepackage[margin=19mm]{geometry}
\usepackage{amsmath,amssymb,mathtools}
\usepackage[scr=rsfs,cal=euler]{mathalfa}
\usepackage{xfrac}
\usepackage[unicode,hidelinks,bookmarks=false]{hyperref}
\newcommand{\ie}{i.e.\ }
\newcommand{\cf}{cf.\ }
\newcommand{\resp}{respectively\ }
\newcommand{\Subsection}{\subsection}
\providecommand{\nobreakdash}{\nobreak\hbox{-}}
\setlength{\emergencystretch}{2em}
\multlinegap0pt
\usepackage{bm}
\usepackage{bbm}
\usepackage{dsfont}
\usepackage{xspace}
\usepackage{cancel}
\usepackage{mathtools}
\usepackage{amsthm}
\makeatletter
\@ifundefined{thm}{\newtheorem{thm}{Thm}}{}
\@ifundefined{lem}{\newtheorem{lem}[thm]{Lemma}}{}
\makeatother
\def\oQ{\overline{Q}}
\title[Decreasing Runs in Quasi-Stirling Permutations of Multisets]{Decreasing Runs in Quasi-Stirling Permutations of Multisets}
\author{Hanqian Fang}
\date{Source: arXiv:2608.09599v1 (CC BY 4.0)}
\begin{document}
\maketitle

\noindent\emph{Source and licence.} Decreasing Runs in Quasi-Stirling Permutations of Multisets. Version arXiv:2608.09599v1 (CC BY 4.0), distributed under the Creative Commons Attribution 4.0 International licence, as recorded in the arXiv licence metadata for this exact version and shown on the abstract page https://arxiv.org/abs/2608.09599v1. Full rights notice: licenses/arxiv-2608.09599-CC-BY-4.0.txt.\par\smallskip

\noindent\textit{Editorial scope.} Counted unit: the Lem whose source label is \texttt{LEM:factor-keep} in the article (source0.tex lines 656--658, source label \texttt{LEM:factor-keep}), delivered together with the article's own complete original proof of it (source0.tex lines 659--678, one \texttt{proof} environment of 20 lines). No mathematics is written, repaired or re-derived: statement and proof are the article's own text line for line. Required context retained uncounted: EQ:form-case2.3 (lines 578--620). Every definition, display and label of those lines is retained unchanged. Omitted from this package and not redistributed: 1--577, 621--655, 679--1200. Nothing inside a retained excerpt is omitted or altered: every retained body is delivered exactly as the article's lines are written, so no omission receipt of any kind accompanies this package. Citations: the retained excerpts carry no citation command at all, so no bibliography entry of the article is transcribed and no reference is redistributed. Reference adaptation: none was needed and none was made; every reference inside the delivered unit resolves inside it, so no unresolved reference or citation is produced. Printed slips kept literal: no printed word, symbol, hypothesis or number of the counted statement or of its proof was altered. Layout and class: the article's own class is \texttt{article} with packages including mathtools, amsmath, amsthm, amsfonts, amssymb, latexsym, mathrsfs, tabularx, diagbox, youngtab, fullpage, color, colortbl, array; the wrapper is the lane's pinned \texttt{amsart} editorial wrapper at 10pt on A4, which restates the theorem-like environments the retained text needs and adds the article's own macro definitions that the retained text needs (\texttt{\textbackslash oQ}), with extra wrapper packages (bm, bbm, dsfont, xspace, cancel, mathtools). The delivered numbering is the unit's own. Assets: no retained span contains a figure environment, a graphics-inclusion command, a PDF-page inclusion, a table or a TikZ picture; no image file is copied into this package, the package declares no asset dependency and no image file is redistributed with it. Licence: version arXiv:2608.09599v1 of this article is distributed under the Creative Commons Attribution 4.0 International licence, as recorded in the arXiv licence metadata for this exact version and shown on the abstract page https://arxiv.org/abs/2608.09599v1. A full rights notice is delivered at licenses/arxiv-2608.09599-CC-BY-4.0.txt. Characters: every retained excerpt is pure ASCII, so the delivered engine, XeLaTeX with its default Latin Modern text font, reports no missing character.
\begin{gather}
\pi=\cdots \underbrace{\rho_1
	\overset{\overset{\text{1st}}{\downarrow}}{i}
	\cdots
\rho_2
\mkern-17mu \overset{\overset{(k_i-1)\text{-th}}{\downarrow}}{i}\mkern-17mu
\cdots}_{\alpha_1}
\underbrace{\rho_3
	\overset{\overset{\text{1st}}{\downarrow}}{(i-1)}
	\cdots
\rho_4
\overset{\overset{k_{i-1}\text{-th}}{\downarrow}}{(i-1)}
\cdots}_{\alpha_2}
\rho_5
\mkern-5mu
\overset{\overset{k_i\text{-th}}{\downarrow}}{i}
\mkern-5mu
\cdots\nonumber\\
%
\rotatebox[origin=c]{-90}{\ $\mapsto$\ }\nonumber\\
%
\Phi(\pi)=\cdots \underbrace{\rho_3
	\overset{\overset{\text{1st}}{\downarrow}}{(i-1)}
	\cdots
\rho_4
\overset{\overset{k_{i-1}\text{-th}}{\downarrow}}{(i-1)}
\cdots
}_{\alpha_2}
\underbrace{\rho_1
	\overset{\overset{\text{1st}}{\downarrow}}{i}
	\cdots
\rho_2
\mkern-17mu
\overset{\overset{(k_{i}-1)\text{-th}}{\downarrow}}{i}
\mkern-17mu
\cdots
}_{\alpha_1}
\rho_5
\mkern-3mu
\overset{\overset{(k_{i-1}+1)\text{-th}}{\downarrow}}{(i-1)}
\mkern-3mu
\cdots\label{EQ:form-case2.3}
\end{gather}

\begin{lem}\label{LEM:factor-keep}
	For any $\pi\in\oQ_M$, $\theta(\pi)$ contains the factor $i(i-1)$ if and only if $\theta(\Phi(\pi))$ does.
\end{lem}

\begin{proof}
	We proceed the proof by cases.
	
	In \textbf{Case 1}, the sequence $S(\pi,(i,i-1))$ (equivalently, $S(\theta(\pi),(i,i-1))$) is necessarily of one of the following two forms:
	\begin{align*}%\label{EQ:form-case1}
	i\cdots i\,(i-1)\cdots(i-1)\,i\cdots i
	\quad\text{or}\quad
	\underbrace{(i-1)\cdots(i-1)}_{\text{may be empty}}\,i\cdots i\,(i-1)\cdots(i-1),
	\end{align*}
	corresponding to {\bf Case 1.1} and {\bf Case 1.2} respectively. Since $k_i>1$, the map's replacement of $i$'s by $i-1$'s leaves the unique factor $i(i-1)$ in $\theta(\pi)$ intact. Thus $\theta(\Phi(\pi))$ retains it as well.
	
	In \textbf{Case 2}, it suffices to show that $\theta(\Phi(\pi))$ has no factor $i(i-1)$. We first consider \textbf{Case 2.2} and \textbf{Case 2.3}. For \textbf{Case 2.2}, it is enough to note that there must be some letter smaller than $i-1$ between the last $i-1$ and the last $i$ in $\pi$. This follows immediately from the assumption that $i-1$ does not occur in $C(\pi,i,k_i-1)$. In \textbf{Case 2.3}, it is clear from~\eqref{EQ:form-case2.3} that $\theta(\pi)$ contains the factor $i(i-1)$ if and only if no letter smaller than $i-1$ lies between the $(k_i-1)$-th $i$ and the first $i-1$ in $\pi$; this condition is equivalent to the presence of the factor $i(i-1)$ in $\theta(\Phi(\pi))$.
	
	Now it remains to prove the lemma for \textbf{Case 2.1}. In this case, no $i-1$ occurs in the subsequence $\pi_{j_1}\pi_{j_1+1}\cdots\pi_{j_2}$, hence none occurs in $\alpha$. Thus a factor $i(i-1)$ in $\theta(\Phi(\pi))$ can only arise from:
	\begin{enumerate}
		\item the adjacent pair formed by $\pi_{\phi(j_3)}$ and the $i-1$ inside $\alpha$ after the replacement, which would require $\pi_{\phi(j_3)}=i$;
		\item the new adjacency created by removing $\alpha$, which would require $\pi_{j_4}=i-1$, where $j_4=\min\{j>\phi(j_2): \pi_j\le i\}$.
	\end{enumerate}
	The first possibility corresponds to the factor $\pi_{\phi(j_3)}\pi_{j_3}=i(i-1)$ in $\theta(\pi)$. The second implies $j_4=j_2$ and thus $\pi_{j_4}=i$. Therefore, both are impossible. This completes the proof.
\end{proof}

\end{document}
